
Table \ref{tab:asImemPer04a} shows the input and outputs of the \asimemg kernel.
\begin{table}[H]
\caption{\asimemg Kernel I/O}
\label{tab:asImemPer04a}
\centering
\begin{tabularx}{\textwidth}{|l |l |l |X|}
  \hline
  Pin & Direction & Width & Explanation \\
  \hline
  \hline
  clk\_i & in & 1 & Core clock (\texttt{clk\_div\_s}); macro clock in function mode \\
  \hline
  tck\_i & in & 1 & JTAG clock; macro clock in load mode \\
  \hline
  im\_mode\_i & in & 1 & 1 = load mode (JTAG IR holds \texttt{I\_MEM}), selects \texttt{tck\_i}; see Table~\ref{tab:asImemModes} \\
  \hline
  \asimemAddrPinInt & in & 15 & From program counter (PC).  There is no byte access, so the lower two address bits will not be needed. In case of a program download, this is a part of the JTAG scan-chain (outside the block). \\
  \hline 
  \asimemDatoPinInt & out & 32 & Instruction. Goes to the instruction decoder. \\
  \hline
  \asimemDatiPinInt & in & 32 & Data input port.  From JTAG scan-chain (JTAG data register). This port is used for a program download. \\
  \hline
  \asimemWePinInt & in & 1 & Write enable for data download. From JTAG scan-chain (JTAG data register). This port is used for a program download. \\
  \hline
\end{tabularx}
\end{table}
